\section{Observation geometry and classical conditioning}
\label{sec:conditioning-literature}
The perturbation of an already specified matrix polynomial and the
inversion of a normalized observation are distinct maps. The distinction
can be expressed at the level of their differentials, and identifies
which ingredients of the preceding proofs are classical.

For a regular matrix polynomial $M$, a simple finite eigenvalue $x$,
and nonzero left and right null vectors $y,v$, differentiation gives
\[
 \dot x=-\frac{y^*\dot M(x)v}{y^*M'(x)v}.
\]
Tisseur's coefficient perturbation framework \cite{Tisseur00} computes
condition numbers by maximizing such a functional over specified
coefficient perturbation balls. Anguas--Bueno--Dopico \cite{ABD19}
compare nonhomogeneous and homogeneous eigenvalue condition numbers,
including the treatment of zero and infinite eigenvalues. Accordingly,
one must specify coefficient weights, the eigenvalue metric, and the
admissible perturbations before comparing condition numbers. Our
leading-invertibility hypotheses and fixed finite-root disk keep the
present target in a finite affine chart; a change to a homogeneous
metric is not implicitly made in an observation norm.

At multiple eigenvalues, the Jordan and polynomial spectral background
is supplied by \cite{GLR82,MBO97}. Kressner--Pel\'aez--Moro \cite{KPM09}
study structured H\"older conditioning of multiple eigenvalues and
matrix pencils. In particular, multiplicity exponents and the effect
of restrictions on perturbations belong to an established theory.
The Laurent exclusion argument in Proposition~\ref{prop:general-bound}
and the argument-principle root counts are instances of this classical
machinery. Neither a reciprocal multiplicity exponent nor an expansion
of a triangular inverse, by itself, is a novelty claim here.

The additional object is the pullback from the probability observation.
In the notation of \eqref{eq:normal-tangent}, an observation tangent $Z$
produces
\[
 \dot M_Z=Q_1^{\mathsf T}I_d(qZ+r_ZP)Q_2,
 \qquad
 r_Z=\cN_P^\dagger[-(\Omega\otimes I)(q(T_j)\operatorname{vec}Z_j)_j].
\]
Substituting this map into the simple-root functional describes the
local observation condition. Bounding its two summands independently
by a single unstructured coefficient condition can multiply the
normalization loss by a second matrix loss. Theorem~\ref{thm:joint}
instead performs the scalar gauge first on complete-polynomial
subspaces. Theorem~\ref{thm:flag-gauge} performs it on a common invariant
flag, even when the coefficient algebra is noncommutative or defective.
Only the remaining arbitrary matrix error is then charged the flag-path
resolvent exponent. These are observation-level statements about a
specific structured pullback, not replacements for coefficient
conditioning theory. They do not assert the existence of such a flag
for every matrix-polynomial algebra.

Theorem~\ref{thm:exact-condition} is the norm of the target differential
on the recovered quotient manifold in the specified observation metric.
Its Gram-matrix expression is elementary once that manifold and all
nearby representatives have been identified. Theorem~\ref{thm:local-minimax}
uses the classical local asymptotic normality and minimax framework
\cite{Hajek72,vdV98}; the finite-permutation Gaussian experiment is the
step needed for the unordered target at collisions. The same observable
matrix $V_H$, not the global certificate $\eta_d$, governs both its
metric and statistical assertions.

Finally, semialgebraic projection, choice and Puiseux expansions in
Theorem~\ref{thm:modulus-minimax} and Corollary~\ref{cor:arc-order} are
classical \cite{BCR98,BPR06}. Their application retains all the stochastic
and real-rooted inequalities in the admissible arcs. This is essential:
the positive sum-of-squares identity in
Lemma~\ref{lem:anchored-real-roots} improves an anchored real-rooted
comparison, whereas the two moving alternatives in
Theorem~\ref{thm:multiplicity} still give the slower fixed-neighbourhood
rate. An unstructured complex coefficient ball, a fixed parameter
neighbourhood and a shrinking probability neighbourhood cannot be
substituted for one another in these conclusions.
